% Right-hand operation in an equation transformation chain.
\NewDocumentCommand{\TransformOp}{m}{%
  {\color{black!55}\Big|\;#1}%
}

% Shared geometry for branching and merging equations.
\newlength{\EquationTemplateUpperWidth}
\newlength{\EquationTemplateLowerWidth}
\newlength{\EquationTemplateMaxWidth}
\newlength{\EquationTemplateUpperNoteWidth}
\newlength{\EquationTemplateLowerNoteWidth}
\newlength{\EquationTemplateMaxNoteWidth}
\NewDocumentCommand{\EquationTemplateMeasure}{m m}{%
  \settowidth{\EquationTemplateUpperWidth}{\ensuremath{\displaystyle #1}}%
  \settowidth{\EquationTemplateLowerWidth}{\ensuremath{\displaystyle #2}}%
  \ifdim\EquationTemplateUpperWidth>\EquationTemplateLowerWidth
    \setlength{\EquationTemplateMaxWidth}{\EquationTemplateUpperWidth}%
  \else
    \setlength{\EquationTemplateMaxWidth}{\EquationTemplateLowerWidth}%
  \fi
}
\NewDocumentCommand{\EquationTemplateMeasureNotes}{m m}{%
  \settowidth{\EquationTemplateUpperNoteWidth}{#1}%
  \settowidth{\EquationTemplateLowerNoteWidth}{#2}%
  \ifdim\EquationTemplateUpperNoteWidth>\EquationTemplateLowerNoteWidth
    \setlength{\EquationTemplateMaxNoteWidth}{\EquationTemplateUpperNoteWidth}%
  \else
    \setlength{\EquationTemplateMaxNoteWidth}{\EquationTemplateLowerNoteWidth}%
  \fi
}
\NewDocumentCommand{\EquationTemplateOverlay}{m m m}{%
  \makebox[0pt][l]{%
    \hspace{#2}\raisebox{#1}[0pt][0pt]{#3}}%
}
\NewDocumentCommand{\EquationTemplateRelation}{m m}{%
  \rotatebox{#2}{\ensuremath{\displaystyle #1}}%
}
\NewDocumentCommand{\EquationTemplateEquals}{m}{%
  \EquationTemplateRelation{=}{#1}%
}
\NewDocumentCommand{\EquationTemplateUse}{m}{%
  \ifmmode
    \vcenter{\hbox{#1}}%
  \else
    #1%
  \fi
}
\NewDocumentCommand{\MathHighlight}{m}{%
  \begingroup
  \setlength{\fboxsep}{1.5pt}%
  \colorbox{yellow!65}{\ensuremath{\displaystyle #1}}%
  \endgroup
}

% One expression continued in two alternative directions:
% \EquationBranches[<relation>]{<start>}{<upper continuation>}{<upper note>}
%                              {<lower continuation>}{<lower note>}
\NewDocumentCommand{\EquationBranchesContent}{O{=} m m m m m}{%
  \begingroup
  \EquationTemplateMeasure{#3}{#5}%
  \EquationTemplateMeasureNotes{#4}{#6}%
  \ensuremath{\displaystyle #2}%
  \hspace{0.15cm}%
  \EquationTemplateOverlay{0.7em}{0pt}{\EquationTemplateRelation{#1}{35}}%
  \EquationTemplateOverlay{1.4em}{0.55cm}{%
    \makebox[\EquationTemplateMaxWidth][l]{\ensuremath{\displaystyle #3}}%
    \hspace{0.55cm}#4}%
  \EquationTemplateOverlay{-0.7em}{0pt}{\EquationTemplateRelation{#1}{-35}}%
  \EquationTemplateOverlay{-1.4em}{0.55cm}{%
    \makebox[\EquationTemplateMaxWidth][l]{\ensuremath{\displaystyle #5}}%
    \hspace{0.55cm}#6}%
  \rule[-1.8em]{0pt}{4.3em}%
  % Reserve the full visible branch width so the complete construction,
  % including both notes, is centered by the surrounding display.
  \hspace{1.10cm}%
  \hspace{\EquationTemplateMaxWidth}%
  \hspace{\EquationTemplateMaxNoteWidth}%
  \endgroup
}

\NewDocumentCommand{\EquationBranches}{O{=} m m m m m}{%
  \EquationTemplateUse{%
    \EquationBranchesContent[#1]{#2}{#3}{#4}{#5}{#6}}%
}

% Two alternative expressions leading to one common result:
% \EquationMerge{<upper note>}{<upper expression>}{<lower note>}
%                {<lower expression>}{<result>}
\NewDocumentCommand{\EquationMergeContent}{m m m m m}{%
  \begingroup
  \EquationTemplateMeasure{#2}{#4}%
  \makebox[\EquationTemplateMaxWidth][l]{%
    \EquationTemplateOverlay{1.4em}{0pt}{%
      \makebox[0pt][r]{#1\hspace{0.55cm}}%
      \makebox[\EquationTemplateMaxWidth][r]{\ensuremath{\displaystyle #2}}}%
    \EquationTemplateOverlay{-1.4em}{0pt}{%
      \makebox[0pt][r]{#3\hspace{0.55cm}}%
      \makebox[\EquationTemplateMaxWidth][r]{\ensuremath{\displaystyle #4}}}%
    \rule[-1.8em]{0pt}{4.3em}}%
  \hspace{0.15cm}%
  \EquationTemplateOverlay{0.7em}{0pt}{\EquationTemplateEquals{-35}}%
  \EquationTemplateOverlay{-0.7em}{0pt}{\EquationTemplateEquals{35}}%
  \hspace{0.55cm}%
  \ensuremath{\displaystyle #5}%
  \endgroup
}

\NewDocumentCommand{\EquationMerge}{m m m m m}{%
  \EquationTemplateUse{%
    \EquationMergeContent{#1}{#2}{#3}{#4}{#5}}%
}

\newcommand{\R}{\mathbb{R}}
\newcommand{\N}{\mathbb{N}}
\newcommand{\Q}{\mathbb{Q}}
\newcommand{\Z}{\mathbb{Z}}
\newcommand{\C}{\mathbb{C}}
% Right-angle mark: a circular arc and a small dot inside its sector.
% The angles describe the visible projection; a spatial right angle need not
% appear as 90 degrees on the page. Radius and style remain figure-specific.
% Arguments: [style]{vertex}{start angle}{end angle}{arc radius}
%            [dot radius][dot distance as a fraction of the arc radius].
\NewDocumentCommand{\RightAngleMark}{O{black!65} m m m m O{0.85pt} O{0.6667}}{%
  \pgfmathsetmacro{\RightAngleMidAngle}{(#3+#4)/2}%
  \pgfmathsetmacro{\RightAngleDotDistance}{#5*(#7)}%
  \draw[#1,line width=0.75pt] ($#2+(#3:#5)$)
    arc[start angle=#3,end angle=#4,radius=#5];%
  \fill[#1] ($#2+(\RightAngleMidAngle:\RightAngleDotDistance)$) circle (#6);%
}
\DeclareMathOperator{\grad}{grad}
\DeclareMathOperator{\arsinh}{arsinh}
\DeclareMathOperator{\arcosh}{arcosh}
\DeclareMathOperator{\artanh}{artanh}
\newcommand{\corresponds}{\mathrel{\widehat{=}}}
